% ch6_c §6.8–6.10 (书页 190–199 / p204–p213)
%--B-TAIL--
%FIG{112}: 电子在一维电场中运动。(a) $k$ 空间中的轨迹。(b) 实空间中的路径。

\section{Zener 击穿：隧穿}

§6.4 中阐述的等效哈密顿量理论，依赖于可以忽略带间跃迁这一假设。这只有在电场不太强时才成立；我们必须能够以不同能带的 Wannier 函数相互正交为理由，略去如下形式的矩阵元
\begin{equation*}
\mathcal{U}_{nn'}(l,l)=\int a^{*}_{n'}(\mathbf{r}-l)\,\mathcal{U}(\mathbf{r})\,a_{n}(\mathbf{r}-l)\,\mathrm{d}\mathbf{r}, \tag{6.50}
\end{equation*}
但如果 $\mathcal{U}(\mathbf{r})$ 随 $\mathbf{r}$ 快速变化，矩阵元 (6.50) 就未必定会为零。从纯数学的观点看，显然带间跃迁的概率将取决于 $\nabla\mathcal{U}$ 的大小。

我们想必能够由 §6.4 的表述出发来计算这个概率。但换一个更“物理的”角度来看待这一现象要容易些。让我们先考虑一个电子在弱电场中会遇到什么情况。如果电子是自由的，它就会被不断地加速，速度无限增大，一路前进，直到被散射，或者跑出样品。

但假定电子处在单一的一条简单能带中——我们可以设想一个一维模型，电场沿 $x$ 轴。代表点便稳稳地前进，从 $O$ 到 $A$。这与 $A'$ 等价——所以我们可以让它突然跳回那里——或者，采用 §3.3 的重复区图式，我们可以看到它沿周期曲线 $OABC$ 等等上下往复运动。

%FIG{113}: 电子不是继续走向 $B$，而是可能“跳”过带隙到达 $A''$。

在实空间中，电子从 $O$ 出发被加速，但当它接近 $A$ 时开始减慢，然后反向运动，直到它到达 $B$，再次被加速，如此往复。在 $A$ 处发生 Bragg 反射；电子不能进入 $A$ 以外的区域，因为它的能量将落入能隙之中。电子被限制在 $x$ 轴的一个有限范围内。

但现在假定还有另一条能带，其最低点 $A''$ 位于 $A$ 之上能量 $\mathcal{E}_{\mathrm{gap}}$ 处。电子总有这种可能：它从电场获得太多的能量，以至于从 $A$ 跃迁到 $A''$（译注：原文印作 to $A$ from $A''$，当为排印倒置）。

为看出这确实是可能的，设电场为恒定的。我们可以沿 $x$ 轴画出电子从电场获得的能量 $eEx$。设电子已到达较低能带顶部的 $A$ 点。那么通常它将被反射回去。但如果它能再走过一段距离
\begin{equation*}
d=\frac{\mathcal{E}_{\mathrm{gap}}}{eE}, \tag{6.51}
\end{equation*}
到达 $A''$ 点，它就已从电场获得了足以越过带隙的能量（图 114）。

%FIG{114}: 被强电场倾斜的能带。

%FIG{115}: 穿过势垒的隧穿。

这本质上是一个\emph{隧穿}（tunnelling）问题。电子必须穿透一个禁戒区域。在传统的隧穿问题中，势垒是这样一个区域：势能 $\mathcal{V}$ 大于电子在自由空间中的总能量 $\mathcal{E}$，因而在势垒内部电子的动能将是负的。从量子力学上看，这意味着电子波函数获得一个实的指数因子
\begin{equation*}
\psi=\psi_{0}e^{-\beta x}, \tag{6.52}
\end{equation*}
其中
\begin{equation*}
\beta^{2}=\mathcal{V}-\mathcal{E}. \tag{6.53}
\end{equation*}
如果 $\mathcal{V}$ 是 $x$ 的函数，我们可以用 W.K.B. 方法导出跃迁概率，即透射与反射之比，其形式为
\begin{equation*}
P=\exp\Biggl\{-2\int_{x_{1}}^{x_{2}}\beta(x)\,\mathrm{d}x\Biggr\}, \tag{6.54}
\end{equation*}
其中 $x_{1}$ 和 $x_{2}$ 是势垒区域的边界，因子 2 来自把波函数平方以得到概率密度。

要把这一论证用到我们眼下的问题上，我们需要考虑能隙中电子波函数的性质。这是我们此前未曾考虑过的一点。偏微分方程
\begin{equation*}
-\frac{\hbar^{2}}{2m}\nabla^{2}\psi+\bigl\{\mathcal{V}(\mathbf{r})-\mathcal{E}\bigr\}\psi=0 \tag{6.55}
\end{equation*}
对任何 $\mathcal{E}$ 值都有解，即使当 $\mathcal{V}(\mathbf{r})$ 是满足
\begin{equation*}
\mathcal{V}(\mathbf{r}+l)=\mathcal{V}(\mathbf{r}), \tag{6.56}
\end{equation*}
的周期势时也是如此；只不过这些解不满足 Bloch 条件，因而不能作为本征函数去配合周期边界条件。

不过，我们并没有理由不采用 §3.2 中那样的势的 Fourier 分析，并尝试如下形式的解
\begin{equation*}
\psi=\sum_{\mathbf{g}}\alpha_{\mathbf{k}-\mathbf{g}}\,e^{i(\mathbf{k}-\mathbf{g})\cdot\mathbf{r}}, \tag{6.57}
\end{equation*}
只是让 $\mathbf{k}$ 的分量取复数值。我们将得到与 (3.14) 同样的系数线性方程组——只是必须把
\begin{equation*}
\mathcal{E}^{0}_{\mathbf{k}-\mathbf{g}}\equiv\frac{\hbar^{2}}{2m}(\mathbf{k}-\mathbf{g})^{2} \tag{6.58}
\end{equation*}
显式地写出来，因为它严格说来已不再是“未微扰能量”，而不过是 $(-\hbar^{2}/2m)\nabla^{2}$ 作用于 (6.57) 中各项的结果。

让我们考察一个一维模型中能隙附近的区域，即 $\mathbf{k}$ 位于与波矢
\begin{equation*}
\tfrac{1}{2}G=\frac{\pi}{a} \tag{6.59}
\end{equation*}
相应的区边界邻域的情形，此处晶面与电场垂直，其间距为 $a$。对自由电子，能量本应是
\begin{equation*}
\mathcal{E}^{0}=\frac{\hbar^{2}}{2m}\bigl(\tfrac{1}{2}G\bigr)^{2} \tag{6.60}
\end{equation*}
但正如 (3.16) 中那样，由于存在微扰 $\mathcal{V}_{G}$，我们必须求解行列式方程
\begin{equation*}
\biggl\{\frac{\hbar^{2}}{2m}k^{2}-\mathcal{E}\biggr\}\biggl\{\frac{\hbar^{2}}{2m}(k-G)^{2}-\mathcal{E}\biggr\}=|\mathcal{V}_{G}|^{2} \tag{6.61}
\end{equation*}
正如我们在 (3.19) 和 (3.20) 中看到的那样，若 $k$ 为实数，则此方程在能隙内没有 $\mathcal{E}$ 的根，也就是说，在
\begin{equation*}
\mathcal{E}^{0}-|\mathcal{V}_{G}|<\mathcal{E}<\mathcal{E}^{0}+|\mathcal{V}_{G}|. \tag{6.62}
\end{equation*}
的范围内无根。

但现在我们把 $\mathcal{E}$ 当作一个任意参量，转而求解 $k$。若写成
\begin{equation*}
k=\tfrac{1}{2}G+\kappa,\qquad\mathcal{E}=\mathcal{E}^{0}+\epsilon, \tag{6.63}
\end{equation*}
其中 $\kappa$ 和 $\epsilon$ 设为小量，则我们求得一个近似解
\begin{equation*}
\kappa^{2}\approx\frac{2m}{\hbar^{2}}\biggl[\frac{\epsilon^{2}-|\mathcal{V}_{G}|^{2}}{4\mathcal{E}^{0}}\biggr]. \tag{6.64}
\end{equation*}

显而易见，若 $|\epsilon|<|\mathcal{V}_{G}|$，则 $\kappa$ 将是纯虚数——也就是说，按 (6.62) 与 (6.63)，\emph{能量处在能隙内的电子，其波函数含有一个实的指数因子}。这完全是 Schrödinger 方程的一个正当的解——但通常在物理上无关紧要，因为在大晶体整个宽度上无法使它配合周期边界条件或其他简单边界条件。电子根本无法在固体中传播任何距离。周期系统中波函数的这种普遍性质，已经在频率位于声子谱之外的晶格振动的情形（§2.12）中指出过。

回到我们的隧穿问题，我们看到：令 (6.64) 中系数 $\beta$ 等于 $\kappa$ 的虚部，便可以用 (6.54) 来计算穿过禁戒区的透射。如果电场使电子能量在跨越禁戒区时线性增大，即 $\epsilon$ 正比于离开势垒中心（势垒总宽度为 $d$）的距离 $x$，那么我们有
\begin{equation*}
\beta(x)=\frac{|\mathcal{V}_{G}|}{2\mathcal{E}^{0}}\bigl(\tfrac{1}{2}G\bigr)\biggl(1-\frac{4x^{2}}{d^{2}}\biggr)^{\frac{1}{2}},
\end{equation*}
积分后得
\begin{align*}
P&=\exp-\biggl\{\frac{\pi d\,|\mathcal{V}_{G}|\,(\tfrac{1}{2}G)}{2\mathcal{E}^{0}}\biggr\}\\
&=\exp-\biggl\{\frac{\pi^{2}}{4}\,\frac{\mathcal{E}_{\mathrm{gap}}^{2}}{\mathcal{E}^{0}eEa}\biggr\}, \tag{6.65}
\end{align*}
这里用到了 (6.51)、(6.59)，以及 $2|\mathcal{V}_{G}|$ 即能隙宽度这一事实。

这个公式表明，如果电场 $E$ 强到电子在走过一个晶格间距 $a$ 内就能获得带隙能量的 $\mathcal{E}_{\mathrm{gap}}/\mathcal{E}^{0}$ 份额，就应当发生 Zener 击穿。然而，这一效应在实践中能否观察到尚不确定，因为它会被其他更复杂的现象所掩盖，例如\emph{雪崩击穿}（avalanche breakdown），即在由高能电子或空穴组成的电流中通过使杂质能级电离而令载流子倍增。

%FIG{116}: 隧道二极管特性：(a) 近零偏压；(b) 隧穿入带隙；(c) 声子辅助隧穿；(d) 少数载流子注入。

一个非常类似的效应出现在\emph{隧道二极管}或 \emph{Esaki 二极管}（Esaki diode）中：在分隔重掺杂 \emph{p} 型与 \emph{n} 型材料的一层薄绝缘膜两端加上外电场。隧穿速率主要由势垒的厚度和性质决定，但也受载流子所要进入的能带中可资利用的态密度的调制——而态密度对于电子被外加电场携带着越过禁戒区时所到达的确切能量颇为敏感。

例如在零偏压下（图 116(a)），势垒两侧的费米能级必须相同（§4.6），于是一个小的“正向”电压会使 \emph{n} 型区的电子隧穿过来，与 \emph{p} 型材料中的空穴相遇而复合，从而使电流得以流动。但势垒两侧更大的电压差会把载流子从一侧驱入另一侧的带隙之中，在那里它们无法传播（图 116(b)）。于是电流下降——但不降至零，因为载流子可以通过与声子的相互作用损失能量（\emph{声子辅助隧穿}（phonon-assisted tunnelling），图 116(c)）。与光子类似的一个效应在光学上表现为 \emph{Franz--Keldysh 效应}（Franz--Keldysh effect）。当然，在更大的电压下，隧穿之后态又变得可资利用了（图 116(d)），于是随着\emph{少数载流子}（minority carriers）的注入——即电子注入 \emph{p} 型材料、空穴注入 \emph{n} 型材料——电流再度上升。

超导体之间的结的行为与此类似，而且还表现出\emph{相干隧穿}（coherent tunnelling）现象，这将在 §11.10 中讨论。

\section{表面上的电子}

第 3 章中讨论的电子波函数被假定满足循环边界条件 (1.76)，仿佛深处大晶体内部。那么在表面附近会发生什么呢？

首先，我们需要对表面各层中原子、离子和电子密度的分布建立某种模型。我们一般可以假定势有一个抬升——即\emph{功函数}（work function）（图 117），说的是把一个电子从体材料的费米能级带到离表面相当远处所需的能量提升——但要计算所需爬越的势垒的确切形式却极其困难。即使所有化学杂质都已清除——这是一项艰巨的技术任务——仍可能有结构重排和电子密度的复杂移动，其起因是未饱和的化学键、镜像力等等。

%FIG{117}: 表面处的功函数势垒。

然而，容易证明，普通 Bloch 态的分布几乎不受表面条件的影响。一个在晶体外呈指数衰减的波函数，可以同一个本来就近乎简并的内部波的组合相匹配，能量只需移动一个无穷小量。匹配的确切细节将依赖于势垒的精确形状，但这对体材料的能带结构等没有影响。

但要考虑能量处在带隙内的电子。如上一节所示，晶体\emph{内部} Schrödinger 方程的解同样由距离的指数函数 (6.52) 主导。设此函数是向晶体内部衰减的，然后我们成功地把它在振幅和梯度上同一个简单指数函数匹配起来——后者描写一个试图穿过表面势垒逃逸出去的电子。这样我们就构造出了一个满足有限性等一般条件的本征态，但它定域在晶体表面处（图 118）。

%FIG{118}: 表面态。

设在某一维模型的一个带隙中，我们已在能量 $\mathcal{E}_{0}$ 处发现了一个定域态 $\psi_{0}(z)$，这个一维模型代表的或许是当我们沿离开表面的 $z$ 方向逐层前进时的平均原子势或赝势。既然晶体在平行于表面的方向上想必仍是周期的，Bloch 定理 (1.58) 对该平面内的平移就必定仍然成立。因此，一维的定域态必定扩展成整整一条\emph{表面态}（surface states）或 \emph{Tamm 态}（Tamm states）的带，其形式为
\begin{equation*}
\psi_{\mathbf{k}}=\psi_{0}(z)\exp i(k_{x}x+k_{y}y), \tag{6.66}
\end{equation*}
其中 $k_{x}$ 和 $k_{y}$ 是在表面平面内量度的波矢的分量。例如在 N.F.E. 情形，我们应预期这条带的行为有如
\begin{equation*}
\mathcal{E}(\mathbf{k})\sim\mathcal{E}_{0}+(\hbar^{2}/2m)(k_{x}^{2}+k_{y}^{2}) \tag{6.67}
\end{equation*}
一直延伸到 $k_{x}$ 与 $k_{y}$ 的倒格子空间中某种二维区边界的类似物为止。

事实上，电子表面态在原理上与 §2.12 中讨论的晶格振动表面模式完全相同。只要对符号作适当的重新标注，我们甚至可以把 (2.138)–(2.145) 的代数演算诠释为在适当边界条件下描写紧束缚模型中的电子表面态。但这类态能够出现的确切能量依赖于模型的细节。即使在一维情形也很显然：势垒的形状及其相对于晶格中最近“原子”的位置，会对 $\psi$ 的局部曲率有举足轻重的影响，从而决定能否实现良好的匹配。在三维情形，情况更加复杂，因为匹配必须在整个表面上进行，而表面绝不会是电子势的一个简单的平面突变。原则上我们可以讨论这类态在给定带隙中出现的条件，但它们对结构细节、化学杂质和空间电荷场的敏感性，阻碍了对能谱及其他性质的定量计算。尽管理论上不可预测，人们还是发现表面态在实际半导体器件的物理中扮演着重要角色：它们可能造成陷阱、复合中心、短路沟道、附加电容以及其他不良效应。这类态在表面催化以及界面上的其他化学现象中也很重要。

真正的表面态只有在功函数势垒之下存在带隙时才能出现。在更高的能量处，与之相应的\emph{倏逝表面波}（evanescent surface waves）在时间上只是准驻定的，类似于共振（参见 §2.12）。这时的典型表面现象便是电子的发射——热发射或光激活发射（见 §8.5）——以及其逆过程：电子从外部射向表面。

高能电子穿过薄膜的衍射已在 §2.6、§3.3 中讨论过。\emph{低能电子衍射}（Low Energy Electron Diffraction, L.E.E.D.）现象则要难作定量解释得多。对于被表面反射回来的电子，晶体动量的法向分量不再是好量子数。衍射峰的基本条件是 (2.79) 的二维类似式，即
\begin{equation*}
\mathbf{k}'_{\parallel}=\mathbf{k}_{\parallel}+\mathbf{g}_{\parallel}, \tag{6.68}
\end{equation*}
其中初始波矢、末态波矢和倒格矢都用它们在平行于表面的平面——即 (6.66) 的 $(x,y)$ 平面——内的投影来表示。反射束的法向分量总可以调节得满足能量守恒条件 $|\mathbf{k}'|^{2}=|\mathbf{k}|^{2}$，因此无论入射束的能量或方向如何，衍射都是允许的。

但要计算衍射强度，我们需要把入射波和反射波同晶体内部 Schrödinger 方程全部解的一个组合相匹配，也就是同固体的\emph{复能带结构}（complex band structure）的所有分支相匹配。不同衍射束的相对强度将依赖于这些内部波函数中相应平面波的相对比例。例如，最强的峰必定出现在入射电子试图以普通 Bloch 态带隙中的能量进入晶体之时——也就是说，接近完整的三维 Bragg--von Laue 衍射条件 (2.79) 处。但在这一区域还可能出现对能量敏感的“共振”效应，此时倏逝表面波可以把电子“导入”不同的衍射束。杂质散射、吸收以及其他非弹性效应，是解释 L.E.E.D. 观测时额外的复杂因素。

\section{电子被杂质散射}

在 §5.2、§5.3 中已经证明，金属中带电杂质的势会受到屏蔽，取 (5.29) 的形式，
\begin{equation*}
\mathcal{U}(\mathbf{r})=\frac{Ze^{2}}{r}\,e^{-\lambda r}, \tag{6.69}
\end{equation*}
其中 $\lambda$ 由 (5.22) 给出。用初等的含时微扰理论容易证明，这个势给出微分散射截面
\begin{equation*}
\sigma(\theta)=\biggl(\frac{2mZe^{2}}{\hbar^{2}}\biggr)^{2}\frac{1}{(K^{2}+\lambda^{2})^{2}}, \tag{6.70}
\end{equation*}
其中 $K=2k_{F}\sin\tfrac{1}{2}\theta$。

对于在半导体中被散射的载流子，类似的结果应当同样成立：此时 $\lambda$ 由 (5.26) 给出，用 $m^{*}$ 代替 $m$，并再乘上一个因子 $1/\epsilon^{2}$ 以顾及介质的介电常数。这不过是等效哈密顿量原理的一次直接应用。

然而，对这些公式不应过分拘泥于字面。它们在求解散射问题时假定了 Born 近似；这对于处于费米能的电子被深势 (6.69) 散射的情形并不很好。更好的做法是使用分波法，
